\subsection{Annihilators and Spectral Control}
\begin{corollary} \label{cor:operator-characteristic-multiple-minimal-polynomial} \hlt{$q_T$ Is a Multiple of $\mu_T$}\\
Let $\mathcal{V}\neq\{\mathbf{0}\}$ be finite-dimensional over $\mathbb{C}$ and let $T\in\LL(\mathcal{V})$.\\
Then the characteristic polynomial $q_T$ is a polynomial multiple of the minimal polynomial $\mu_T$.
\begin{proof}
The polynomial $q_T$ is nonzero because $\deg q_T=\dim(\mathcal{V})\geq1$ by Proposition \ref{prop:operator-characteristic-polynomial-degree-roots}(i).\\
The identity $q_T(T)=0$ holds by Theorem \ref{thm:cayley-hamilton-operator}. Hence $\mu_T\mid q_T$ by Theorem \ref{thm:operator-annihilator-iff-multiple-minimal-polynomial}.\\
\hlt{Structural Commentary}: Cayley--Hamilton turns the characteristic polynomial into an annihilator.\\
\hlt{Key Insight}: The minimal polynomial divides every annihilator, hence divides $q_T$.
\end{proof}
\end{corollary}

\begin{theorem} \label{thm:operator-minimal-polynomial-roots-eigenvalues} \hlt{Roots of $\mu_T$ Are Eigenvalues}\\
Let $\mathcal{V}\neq\{\mathbf{0}\}$ be finite-dimensional over $\mathbb{F}$ and let $T\in\LL(\mathcal{V})$. Then $(\forall\,\lambda\in\mathbb{F})\,[\mu_T(\lambda)=0\Longleftrightarrow\lambda\in\operatorname{spec}(T)]$.
\begin{proof}
\begin{enumerate}[label=\roman*.]
\setlength{\itemsep}{0pt}
\item[$(\Rightarrow)$] Fix $\lambda\in\mathbb{F}$ with $\mu_T(\lambda)=0$. Termwise, $\mu_T(z)-\mu_T(\lambda)\in(z-\lambda)\mathbb{F}[z]$ by $z^r-\lambda^r=(z-\lambda)\sum\limits_{s=0}^{r-1}z^{r-1-s}\lambda^s$.\\
Thus $\mu_T(z)=(z-\lambda)g(z)$ for a monic $g$ with $\deg g=\deg\mu_T-1$.\\
Hence $g(T)\neq0$ by Definition \ref{def:operator-minimal-polynomial}; choose $\mathbf{v}\in\mathcal{V}$ with $\mathbf{u}\defeq g(T)\mathbf{v}\neq\mathbf{0}$.\\
Then $(T-\lambda I_{\mathcal{V}})\mathbf{u}=\mu_T(T)\mathbf{v}=\mathbf{0}$ by Proposition \ref{prop:operator-polynomial-products} and Definition \ref{def:operator-minimal-polynomial}. Thus $\lambda\in\operatorname{spec}(T)$ by Definition \ref{def:spectrum-and-geometric-multiplicity}.
\item[$(\Leftarrow)$] Fix $\lambda\in\mathbb{F}$ and suppose $\lambda\in\operatorname{spec}(T)$. By Definition \ref{def:spectrum-and-geometric-multiplicity}, choose $\mathbf{v}\neq\mathbf{0}$ satisfying $T\mathbf{v}=\lambda\mathbf{v}$.\\
Then $\mathbf{0}=\mu_T(T)\mathbf{v}=\mu_T(\lambda)\mathbf{v}$ by Definition \ref{def:operator-minimal-polynomial} and Proposition \ref{prop:polynomial-eigenvector-evaluation}. Since $\mathbf{v}\neq\mathbf{0}$, $\mu_T(\lambda)=0$.
\end{enumerate}
\hlt{Structural Commentary}: Root factorization supplies the forward eigenvector; evaluation proves the converse.\\
\hlt{Key Insight}: Minimality keeps the quotient operator nonzero.
\end{proof}
\end{theorem}

\begin{theorem} \label{thm:minimal-polynomial-direct-sum-disjoint-spectra} \hlt{Minimal Polynomial of a Spectrally Disjoint Direct Sum}\\
Let $d,n_1,\ldots,n_d\geq1$, let $A_i\in\operatorname{Mat}(n_i,n_i,\mathbb{C})$ for $i\in\{1,\ldots,d\}$, and set $A\defeq A_1\oplus\cdots\oplus A_d$.\\
If the spectra of $A_1,\ldots,A_d$ are pairwise disjoint, then $\mu_A=\mu_{A_1}\mu_{A_2}\cdots\mu_{A_d}$.
\begin{proof}
Set $f\defeq\prod_{i=1}^{d}\mu_{A_i}$. The polynomial $f$ is monic and hence nonzero.
\begin{enumerate}[label=(\roman*)]
\setlength{\itemsep}{0pt}
\item \hlt{Block Annihilation}: Since $\mu_{A_i}\mid f$, $(\forall\,i\in\{1,\ldots,d\})\,f(A_i)=0$ by Theorem \ref{thm:operator-annihilator-iff-multiple-minimal-polynomial} and Definition \ref{def:minimal-polynomial-matrix}.\\
Hence $f(A)=f(A_1)\oplus\cdots\oplus f(A_d)=0$ by Proposition \ref{prop:direct-sum-matrix-algebra}(iv), and $\mu_A\mid f$ by the same criterion.
\item \hlt{Block Restriction}: The identity $\mu_A(A)=\mu_A(A_1)\oplus\cdots\oplus\mu_A(A_d)=0$ holds by Definition \ref{def:minimal-polynomial-matrix} and Proposition \ref{prop:direct-sum-matrix-algebra}(iv).\\
Thus $(\forall\,i\in\{1,\ldots,d\})\,\mu_{A_i}\mid\mu_A$ by Theorem \ref{thm:operator-annihilator-iff-multiple-minimal-polynomial} and Definition \ref{def:minimal-polynomial-matrix}.
\item \hlt{Coprime Factors}: The roots of $\mu_{A_i}$ are $\operatorname{spec}(A_i)$ by Theorem \ref{thm:operator-minimal-polynomial-roots-eigenvalues}, so the polynomials $\mu_{A_1},\ldots,\mu_{A_d}$ are pairwise coprime.\\
Hence $f\mid\mu_A$ by part (ii) and factorization over $\mathbb{C}$.
\item \hlt{Conclusion}: By parts (i) and (iii), $\mu_A\mid f$ and $f\mid\mu_A$. Monicity forces $\mu_A=f$.
\end{enumerate}
\hlt{Structural Commentary}: Block evaluation establishes $\mu_A\mid f$; spectral disjointness establishes $f\mid\mu_A$.\\
\hlt{Key Insight}: Every spectral block contributes a distinct collection of factors to $\mu_A$.
\end{proof}
\end{theorem}

\begin{corollary} \label{cor:operator-minimal-characteristic-same-roots} \hlt{Minimal and Characteristic Polynomials Have the Same Roots}\\
Let $\mathcal{V}\neq\{\mathbf{0}\}$ be finite-dimensional over $\mathbb{C}$ and let $T\in\LL(\mathcal{V})$. Then $\{\lambda\in\mathbb{C}:\mu_T(\lambda)=0\} = \{\lambda\in\mathbb{C}:q_T(\lambda)=0\} = \operatorname{spec}(T)$.
\begin{proof}
By Theorem \ref{thm:operator-minimal-polynomial-roots-eigenvalues}, $\{\lambda\in\mathbb{C}:\mu_T(\lambda)=0\}=\operatorname{spec}(T)$. By Proposition \ref{prop:operator-characteristic-polynomial-degree-roots}, $\{\lambda\in\mathbb{C}:q_T(\lambda)=0\}=\operatorname{spec}(T)$.\\
\hlt{Structural Commentary}: Different multiplicities, same spectral scalars.\\
\hlt{Key Insight}: Minimal polynomial roots and characteristic polynomial roots are both just eigenvalues.
\end{proof}
\end{corollary}

\begin{perspective} \label{per:characteristic-versus-minimal-polynomial} \hlt{Characteristic vs. Minimal Polynomial}\\
By Corollary \ref{cor:operator-minimal-characteristic-same-roots}, $q_T$ and $\mu_T$ have the same roots.\\
They encode multiplicity differently: $q_T$ records algebraic multiplicity, whereas $\mu_T$ uses the least exponents that annihilate $T$.
\end{perspective}
